\documentclass[10pt,a4paper]{amsart}
\usepackage[margin=19mm]{geometry}
\usepackage{amsmath,amssymb,mathtools}
\usepackage[scr=rsfs,cal=euler]{mathalfa}
\usepackage{xfrac}
\usepackage[unicode,hidelinks,bookmarks=false]{hyperref}
\newcommand{\ie}{i.e.\ }
\newcommand{\cf}{cf.\ }
\newcommand{\resp}{respectively\ }
\newcommand{\Subsection}{\subsection}
\providecommand{\nobreakdash}{\nobreak\hbox{-}}
\setlength{\emergencystretch}{2em}
\multlinegap0pt
\usepackage{mathtools}
\usepackage{amssymb}
\usepackage[shortlabels]{enumitem}
\usepackage{xcolor}
\newtheorem{definition}{Definition}
\newtheorem{lemma}{Lemma}
\title{$J$-EQUATIONS AND DEFORMED HERMITIAN-YANG-MILLS EQUATIONS ON HOLOMORPHIC SUBMERSIONS}
\author{Rei Murakami}
\date{Source: arXiv:2208.08576v4 (CC BY 4.0)}
\begin{document}
\maketitle

\noindent\emph{Source and licence.} $J$-EQUATIONS AND DEFORMED HERMITIAN-YANG-MILLS EQUATIONS ON HOLOMORPHIC SUBMERSIONS. Version arXiv:2208.08576v4 (CC BY 4.0), distributed by arXiv under the Creative Commons Attribution 4.0 International licence, http://creativecommons.org/licenses/by/4.0/, as recorded in the arXiv licence metadata for this exact version and shown on the abstract page https://arxiv.org/abs/2208.08576v4. Full licence notice: \texttt{licenses/arxiv-2208.08576-CC-BY-4.0.txt}.\par\smallskip

\noindent\textit{Editorial scope.} Counted unit: the article's lemma lemma (source0.tex lines 1019--1025). The article's complete original proof of it is delivered with it (source0.tex lines 1027--1062, one \texttt{proof} environment of 36 lines). No mathematics is written, repaired or re-derived: the statement and the proof are the article's own lines, in the article's own notation, and no retained line is substituted or re-flowed.  Retained uncounted as the source-local context the proof needs: lines 187--190.  Nothing was omitted from any retained span, so the package carries no omission record.  No retained line contains a citation, so the wrapper carries no bibliography and no bibliography entry.  No retained line contains a figure environment, an image-inclusion command or drawing code, so no image is delivered and the package declares no asset dependency.  Editorial formatting only, in addition: an AMS \texttt{amsart} preamble that restates the article's own declarations and macros used by the retained lines, namely the environments \texttt{definition}, \texttt{lemma} on separate counters, together with the standard packages those lines need.  The retained environments print in this document as Definition 1, Lemma 1 in order of appearance, whereas the article prints the same items with the numbering of its own full document.  The complete native source of the article as distributed in the arXiv e-print for this exact version is bundled unchanged as \texttt{sources/129596/source0.tex}.
\begin{definition}\label{def:F-op}
We define an operator $F_{\omega, \chi} : C^{\infty}(X) \rightarrow C^{\infty}(X) $ by 
$$F_{\omega, \chi} (\phi) = {-{ \left( \chi , {\sqrt{-1} \partial \bar{\partial} \phi} \right)_\omega} - \left( \partial \left( \Lambda_\omega \chi \right) , \bar{\partial} \phi \right)_\omega }. $$
\end{definition}

\begin{lemma}
 For $\phi \in C^\infty(X)$, we have
$$\frac{d}{dt}\Big|_{t = 0}\frac{\mathrm{Re}(\omega_{k,r,t\phi}+\sqrt{-1}\chi)^{m+n}}{\mathrm{Im}(\omega_{k,r,t\phi}+\sqrt{-1}\chi)^{m+n}}
= \tilde{F}_{ \left( \omega_X \right)_\mathcal{V}, \chi_\mathcal{V}} \left(\phi \right) + k^{-1}D_1(\phi) + k^{-2}D_2(\phi)+O(k^{-3}),$$
where the operators $D_1$ and $D_2$ satisfy $D_1(\phi)=0$ and $\tilde{\pi}_B\bigl(D_2(\phi)\bigr)=-F_{\omega_B,\chi_B}(\phi)/(m+1)$ for $\phi\in C^\infty(B)$. Here, the operator $F_{\omega_B,\chi_B}$ is given by Definition \ref{def:F-op} and we suppress pullbacks via $\pi$,
so for a function $\phi\in C^\infty(B)$ its pullback to $X$ is also denoted by $\phi$.
\end{lemma}

\begin{proof}
    We have
    \begin{align*}
        &\frac{d}{dt}\Big|_{t = 0}\frac{\mathrm{Re}(\omega_{k,r,t\phi}+\sqrt{-1}\chi)^{m+n}}{\mathrm{Im}(\omega_{k,r,t\phi}+\sqrt{-1}\chi)^{m+n}}\\
        =&(m+n)\frac{\mathrm{Re}(\omega_{k,r}+\sqrt{-1}\chi)^{m+n-1}\wedge\sqrt{-1}\partial\bar{\partial}\phi}{\mathrm{Im}(\omega_{k,r}+\sqrt{-1}\chi)^{m+n}}\\
        &-(m+n)\frac{\mathrm{Re}(\omega_{k,r}+\sqrt{-1}\chi)^{m+n}}{\mathrm{Im}(\omega_{k,r}+\sqrt{-1}\chi)^{m+n}}\frac{\mathrm{Im}(\omega_{k,r}+\sqrt{-1}\chi)^{m+n-1}\wedge\sqrt{-1}\partial\bar{\partial}\phi}{\mathrm{Im}(\omega_{k,r}+\sqrt{-1}\chi)^{m+n}}\\
        =&m\frac{\mathrm{Re}((\omega_X)_\mathcal{V}+\sqrt{-1}\chi_\mathcal{V})^{m-1}\wedge(\sqrt{-1}\partial\bar{\partial}\phi)_\mathcal{V}}{\mathrm{Im}((\omega_X)_\mathcal{V}+\sqrt{-1}\chi_\mathcal{V})^m}\\
        &-m\frac{\mathrm{Re}((\omega_X)_\mathcal{V}+\sqrt{-1}\chi_\mathcal{V})^m}{\mathrm{Im}((\omega_X)_\mathcal
        V+\sqrt{-1}\chi_\mathcal{V})^m}\frac{\mathrm{Im}((\omega_X)_\mathcal{V}+\sqrt{-1}\chi_\mathcal{V})^{m-1}\wedge(\sqrt{-1}\partial\bar{\partial}\phi)_\mathcal{V}}{\mathrm{Im}((\omega_X)_\mathcal{V}+\sqrt{-1}\chi_\mathcal{V})^m}\\
        &+k^{-1}D_1(\phi)+k^{-2}D_2(\phi)+O(k^{-3})\\
        =&\tilde{F}_{(\omega_X)_\mathcal{V},\chi_\mathcal{V}}(\phi)+k^{-1}D_1(\phi)+k^{-2}D_2(\phi)+O(k^{-3}),
    \end{align*}
    where the first equation follows since $\sqrt{-1}\partial\bar{\partial}\phi$ is real.
    For $\phi\in C^\infty(B)$, noting that $\sqrt{-1}\partial\bar{\partial}\phi$ has only a horizontal part, 
    we have
    \begin{align*}
        D_1(\phi)
        =&n\frac{\mathrm{Re}((\omega_X)_\mathcal{V}+\sqrt{-1}\chi_\mathcal{V})^m\wedge\omega_B^{n-1}\wedge\sqrt{-1}\partial\bar{\partial}\phi}{\mathrm{Im}((\omega_X)_\mathcal{V}+\sqrt{-1}\chi_\mathcal{V})^m\wedge\omega_B^n}\\
        &-n\frac{\mathrm{Re}((\omega_X)_\mathcal{V}+\sqrt{-1}\chi_\mathcal{V})^m}{\mathrm{Im}((\omega_X)_\mathcal{V}+\sqrt{-1}\chi_\mathcal{V})^m}\frac{\mathrm{Im}((\omega_X)_\mathcal{V}+\sqrt{-1}\chi_\mathcal{V})^m\wedge\omega_B^{n-1}\wedge\sqrt{-1}\partial\bar{\partial}\phi}{\mathrm{Im}((\omega_X)_\mathcal{V}+\sqrt{-1}\chi_\mathcal{V})^m\wedge\omega_B^n}\\
        =&0
    \end{align*}
    Moreover, we have
    \begin{align*}
        &\tilde{\pi}_BD_2(\phi)\\
        =&\tilde{\pi}_B\Biggl(\frac{n(n-1)}{m+1}\frac{\mathrm{Re}(\omega_X+\sqrt{-1}\partial\bar{\partial}\phi_{2,B}+\sqrt{-1}\chi)^{m+1}\wedge\omega_B^{n-2}\wedge\sqrt{-1}\partial\bar{\partial}\phi}{\mathrm{Im}((\omega_X)_\mathcal{V}+\sqrt{-1}\chi_\mathcal{V})^m\wedge\omega_B^n}\\
        &\qquad-\frac{n(n-1)}{m+1}\frac{\mathrm{Re}((\omega_X)_\mathcal{V}+\sqrt{-1}\chi_\mathcal{V})^m}{\mathrm{Im}((\omega_X)_\mathcal{V}+\sqrt{-1}\chi_\mathcal{V})^m}\frac{\mathrm{Im}(\omega_X+\sqrt{-1}\partial\bar{
        \partial}\phi_{2,B}+\sqrt{-1}\chi)^{m+1}\wedge\omega_B^{n-2}\wedge\sqrt{-1}\partial\bar{\partial}\phi}{\mathrm{Im}((\omega_X)_\mathcal{V}+\sqrt{-1}\chi_\mathcal{V})^m\wedge\omega_B^n}\Biggr)\\
        &+\tilde{\pi}_B\Biggl(nm\frac{\mathrm{Re}((\omega_X)_\mathcal{V}+\sqrt{-1}\chi_\mathcal{V})^{m-1}\wedge\sqrt{-1}\partial\bar{\partial}\phi_{1,\mathcal{V}}\wedge\omega_B^{n-1}\wedge\sqrt{-1}\partial\bar{\partial}\phi}{\mathrm{Im}((\omega_X)_\mathcal{V}+\sqrt{-1}\chi_\mathcal{V})^m\wedge\omega_B^n}\\
        &\qquad-nm\frac{\mathrm{Re}((\omega_X)_\mathcal{V}+\sqrt{-1}\chi_\mathcal{V})^m}{\mathrm{Im}((\omega_X)_\mathcal{V}+\sqrt{-1}\chi_\mathcal{V})^m}\frac{\mathrm{Im}((\omega_X)_\mathcal{V}+\sqrt{-1}\chi_\mathcal{V})^{m-1}\wedge
        \sqrt{-1}\partial\bar{\partial}\phi_{1,\mathcal{V}}\wedge\omega_B^{n-1}\wedge\sqrt{-1}\partial\bar{\partial}\phi}{\mathrm{Im}((\omega_X)_\mathcal{V}+\sqrt{-1}\chi_\mathcal{V})^m\wedge\omega_B^n}\Biggr)\\
        &-\tilde{\pi}_B\Biggl(c_1n\frac{\omega_B^{n-1}\wedge\sqrt{-1}\partial\bar{\partial}\phi}{\omega_B^n}\Biggr)\\
        =-&\frac{n(n-1)}{m+1}\frac{\chi_B\wedge\omega_B^{n-2}\wedge\sqrt{-1}\partial\bar{\partial}\phi}{\omega_B^n}-c_1n\frac{\omega_B^{n-1}\wedge\sqrt{-1}\partial\bar{\partial}\phi}{\omega_B^n}\\
        =-&F_{\omega_B,\chi_B}(\phi)/(m+1).
    \end{align*}
    Thus, we confirm the assertion.
\end{proof}

\end{document}
